\documentclass[10pt,a4paper]{amsart}
\usepackage[margin=19mm]{geometry}
\usepackage{amsmath,amssymb,mathtools}
\usepackage[scr=rsfs,cal=euler]{mathalfa}
\usepackage{xfrac}
\usepackage[unicode,hidelinks,bookmarks=false]{hyperref}
\newcommand{\ie}{i.e.\ }
\newcommand{\cf}{cf.\ }
\newcommand{\resp}{respectively\ }
\newcommand{\Subsection}{\subsection}
\providecommand{\nobreakdash}{\nobreak\hbox{-}}
\setlength{\emergencystretch}{2em}
\multlinegap0pt
\usepackage{bm}
\usepackage{bbm}
\usepackage{dsfont}
\usepackage{xspace}
\usepackage{cancel}
\usepackage{mathtools}
\usepackage{amsthm}
\makeatletter
\@ifundefined{claim}{\newtheorem{claim}{Claim}}{}
\@ifundefined{theorem}{\newtheorem{theorem}{Theorem}}{}
\makeatother
\title[On the Subgroup Distance Problem in Cyclic Permutation Group]{On the Subgroup Distance Problem in Cyclic Permutation Groups}
\author{Andreas Rosowski}
\date{Source: arXiv:2504.06844v2 (CC BY 4.0)}
\begin{document}
\maketitle

\noindent\emph{Source and licence.} On the Subgroup Distance Problem in Cyclic Permutation Groups. Version arXiv:2504.06844v2 (CC BY 4.0), distributed under the Creative Commons Attribution 4.0 International licence, as recorded in the arXiv licence metadata for this exact version and shown on the abstract page https://arxiv.org/abs/2504.06844v2. Full rights notice: licenses/arxiv-2504.06844-CC-BY-4.0.txt.\par\smallskip

\noindent\textit{Editorial scope.} Counted unit: the Theorem whose source label is \texttt{theoremlinftyl1l2} in the article (source0.tex lines 797--799, source label \texttt{theoremlinftyl1l2}), delivered together with the article's own complete original proof of it (source0.tex lines 800--959, one \texttt{proof} environment of 160 lines). No mathematics is written, repaired or re-derived: statement and proof are the article's own text line for line. Required context retained uncounted: none. Every definition, display and label of those lines is retained unchanged. Omitted from this package and not redistributed: 1--796, 960--1241. Nothing inside a retained excerpt is omitted or altered: every retained body is delivered exactly as the article's lines are written, so no omission receipt of any kind accompanies this package. Citations: the retained excerpts carry no citation command at all, so no bibliography entry of the article is transcribed and no reference is redistributed. Reference adaptation: none was needed and none was made; every reference inside the delivered unit resolves inside it, so no unresolved reference or citation is produced. Printed slips kept literal: no printed word, symbol, hypothesis or number of the counted statement or of its proof was altered. Layout and class: the article's own class is \texttt{article} with packages including amssymb, amsmath, amsfonts, amsthm, algorithm, a4wide, hyperref, enumerate, mathtools; the wrapper is the lane's pinned \texttt{amsart} editorial wrapper at 10pt on A4, which restates the theorem-like environments the retained text needs and adds the article's own macro definitions that the retained text needs (none), with extra wrapper packages (bm, bbm, dsfont, xspace, cancel, mathtools). The delivered numbering is the unit's own. Assets: no retained span contains a figure environment, a graphics-inclusion command, a PDF-page inclusion, a table or a TikZ picture; no image file is copied into this package, the package declares no asset dependency and no image file is redistributed with it. Licence: version arXiv:2504.06844v2 of this article is distributed under the Creative Commons Attribution 4.0 International licence, as recorded in the arXiv licence metadata for this exact version and shown on the abstract page https://arxiv.org/abs/2504.06844v2. A full rights notice is delivered at licenses/arxiv-2504.06844-CC-BY-4.0.txt. Characters: every retained excerpt is pure ASCII, so the delivered engine, XeLaTeX with its default Latin Modern text font, reports no missing character.
\begin{theorem}\label{theoremlinftyl1l2}
Let $t \in \mathbb{N}$ be odd and let $0 \leq t_1 < t_2 < t$ be such that $t_1 \not \equiv t_2 \bmod p$ for all primes $p$ with $p \mid t$. Then there is a cycle $\alpha$ of length $t$ and a permutation $\beta$ in which $\beta$ is a product of disjoint $2$-cycles such that $l_\infty(\beta, \alpha^{t_1}) \leq 1$ and $l_\infty(\beta, \alpha^{t_2}) \leq 1$.
\end{theorem}

\begin{proof}
We define
\begin{displaymath}
\omega = t_2-t_1.
\end{displaymath}
Then $\omega$ is a generator of the additive group $(\mathbb{Z}_t,+)$ since $t_1 \not \equiv t_2 \bmod p$ for all primes $p$ with $p \mid t$ and in particular $\omega$ and $t$ are coprime and we can define $0 \leq \psi < t$ as the smallest positive integer satisfying
\begin{displaymath}
\psi \equiv \omega^{-1}(t-t_1) \bmod t
\end{displaymath}
since $\omega^{-1} \bmod t$ exists. For $i=0,\dots,t-1$ we define $0 \leq \omega_i < t$ as the smallest positive integer satisfying $\omega_i \equiv i\omega \bmod t$. Moreover for $i=0,\dots,t-1$ we define $0 \leq \psi_i < t$ as the smallest positive integer satisfying $\psi_i \equiv \psi + i \bmod t$. Now we define the cycle $\alpha = (\alpha_0,\dots,\alpha_{t-1})$ of length $t$ by the following:
\begin{displaymath}
\alpha_{\omega_i} = \begin{cases}
2i+1 & \text{if } 0 \leq i \leq \frac{t-1}{2}\\
2(t-i) & \text{if } \frac{t+1}{2} \leq i \leq t-1.
\end{cases}
\end{displaymath}
For $i=0,\dots,t-1$ we define $0 \leq d_{i,1},d_{i,2} < t$ as the smallest positive integers satisfying
\begin{displaymath}
d_{i,k} \equiv i + t_k \bmod t
\end{displaymath}
for $k \in [1,2]$. Then $\alpha^{t_k}$ maps $\alpha_i$ to $\alpha_{d_{i,k}}$ for all $i \in [0,t-1]$.

\begin{claim}\label{claimreldi1di2}
Let $i \in [0,t-1]$ and let $j \in [0,t-1]$ be such that $\omega_j = d_{i,1}$. If $j=t-1$ then $d_{i,2} = \omega_0$ and if $0 \leq j \leq t-2$ then $d_{i,2} = \omega_{j+1}$.
\end{claim}
We have
\begin{eqnarray*}
& d_{i,1} & \equiv \ \ i + t_1 \bmod t\\
\Leftrightarrow & i & \equiv \ \ d_{i,1} - t_1 \bmod t\\
\Leftrightarrow & d_{i,2} \equiv i + t_2 & \equiv \ \ d_{i,1} - t_1 + t_2 \equiv d_{i,1} + \omega \bmod t\\
\Leftrightarrow & d_{i,2} & \equiv \ \ \omega_j + \omega \bmod t\\
\Leftrightarrow & d_{i,2} & \equiv \ \ j\omega + \omega \bmod t\\
\Leftrightarrow & d_{i,2} & \equiv \ \ (j+1)\omega \bmod t\\
\Leftrightarrow & d_{i,2} & \equiv \ \ \begin{cases}
\omega_0 \bmod t & \text{if } j = t-1\\
\omega_{j+1} \bmod t & \text{if } j \in [0,l-2].
\end{cases}
\end{eqnarray*}
Since $0 \leq d_{i,2} < t$ and $0 \leq \omega_0,\omega_{j+1} < t$ we finally obtain
\begin{displaymath}
d_{i,2} = \begin{cases}
\omega_0 & \text{if } j = t-1\\
\omega_{j+1} & \text{if } j \in [0,t-2].
\end{cases}
\end{displaymath}
\qed

By Claim~\ref{claimreldi1di2} we have $1 \leq |\alpha_{d_{i,1}} - \alpha_{d_{i,2}}| \leq 2$ for all $i \in [0,t-1]$ and if $|\alpha_{d_{i,1}} - \alpha_{d_{i,2}}| = 1$ then either $\alpha_{d_{i,2}} = 1$ or $\alpha_{d_{i,1}} = t$. We define for all $i \in [0,t-1]$
\begin{displaymath}
\alpha_i' = \begin{cases}
\frac{\alpha_{d_{i,1}}+\alpha_{d_{i,2}}}{2} & \text{if } |\alpha_{d_{i,1}} - \alpha_{d_{i,2}}| = 2\\
1 & \text{if } |\alpha_{d_{i,1}} - \alpha_{d_{i,2}}| = 1 \text{ and } \alpha_{d_{i,2}} = 1\\
t & \text{if } |\alpha_{d_{i,1}} - \alpha_{d_{i,2}}| = 1 \text{ and } \alpha_{d_{i,1}} = t.
\end{cases}
\end{displaymath}
Now we define the permutation $\beta$ by the following:
\begin{displaymath}
\beta = \prod_{i=0}^{t-1} \beta_i
\end{displaymath}
in which
\begin{displaymath}
\beta_i = \begin{cases}
(\alpha_i, \alpha_i') & \text{if } \alpha_i < \alpha_i'\\
\mathrm{id} & \text{otherwise}.
\end{cases}
\end{displaymath}

\begin{claim}\label{claimreldomegaidomegai1}
For all $i \in [0,t-1]$ we have $d_{\omega_{\psi_i},1} = \omega_i$.
\end{claim}
We have
\begin{align*}
d_{\omega_{\psi_i},1} &\equiv \omega_{\psi_i} + t_1 \bmod t\\
&\equiv \psi_i \omega + t_1 \bmod t\\
&\equiv ((t-t_1)\omega^{-1}+i)\omega + t_1 \bmod t\\
&\equiv t-t_1+i\omega+t_1 \bmod t\\
&\equiv \omega_i \bmod t.
\end{align*}
Since $0 \leq d_{\omega_{\psi_i},1},\omega_i < t$ we finally obtain $d_{\omega_{\psi_i},1} = \omega_i$.
\qed

\begin{claim}\label{claimrelpsiil1i}
For all $i \in [0,t-1]$ we have $\alpha'_{\omega_{\psi_i}} = \alpha_{\omega_{t-1-i}}$.
\end{claim}
Note that by Claim~\ref{claimreldomegaidomegai1} we have $d_{\omega_{\psi_i},1} = \omega_i$. Then we have by Claim~\ref{claimreldi1di2}
\begin{displaymath}
d_{\omega_{\psi_i},2} = \begin{cases}
\omega_{i+1} & \text{if } i \in [0,t-2]\\
\omega_0 & \text{if } i = t-1.
\end{cases}
\end{displaymath}
Case~1: $i \in [0,\frac{t-3}{2}]$. We have
\begin{align*}
\alpha'_{\omega_{\psi_i}}&= \frac{\alpha_{d_{\omega_{\psi_i},1}} + \alpha_{d_{\omega_{\psi_i},2}}}{2}\\
&= \frac{\alpha_{\omega_i} + \alpha_{\omega_{i+1}}}{2}\\
&= \frac{2i+1 + 2(i+1)+1}{2}\\
&= 2i+2\\
&= 2(t-(t-1-i))\\
&= \alpha_{\omega_{t-1-i}}.
\end{align*}
Case~2: $i = \frac{t-1}{2}$. We have $\alpha_{d_{\omega_{\psi_i},1}} = \alpha_{\omega_i} = t$ and hence
\begin{displaymath}
\alpha'_{\omega_{\psi_i}} = \alpha_{d_{\omega_{\psi_i},1}} = t = \alpha_{\omega_{t-1-\frac{t-1}{2}}}.
\end{displaymath}
Case~3: $i \in [\frac{t+1}{2},t-2]$. We have
\begin{align*}
\alpha'_{\omega_{\psi_i}}&= \frac{\alpha_{d_{\omega_{\psi_i},1}} + \alpha_{d_{\omega_{\psi_i},2}}}{2}\\
&= \frac{\alpha_{\omega_i} + \alpha_{\omega_{i+1}}}{2}\\
&= \frac{2(t-i) + 2(t-(i+1))}{2}\\
&= 2(t-i)-1\\
&= 2(t-1-i)+1\\
&= \alpha_{\omega_{t-1-i}}.
\end{align*}
Case~4: $i=t-1$. We have $\alpha_{d_{\omega_{\psi_i},2}} = \alpha_{\omega_0} = 1$ and hence
\begin{displaymath}
\alpha'_{\omega_{\psi_i}} = \alpha_{d_{\omega_{\psi_i},2}} = 1 = \alpha_{\omega_{t-1-(t-1)}}.
\end{displaymath}
\qed

\begin{claim}\label{claimrelalphaialphaj}
For all $i \in [0,t-1]$ and $j \in [0,t-1]$ we have $\alpha_i' = \alpha_j$ if and only if $\alpha_j' = \alpha_i$.
\end{claim}
Suppose $\alpha_i' = \alpha_j$ and let $0 \leq e < t$ be such that $i = \omega_{\psi_e}$. Note that $e$ exists. Since $\omega$ is a generator there is $c \in [0,t-1]$ such that $i \equiv \omega_c \bmod t$ and we can choose $e \equiv c - \psi \bmod t$. By Claim~\ref{claimrelpsiil1i} we have $\alpha'_{\omega_{\psi_e}} = \alpha_{\omega_{t-1-e}}$ (i.e. $j = \omega_{t-1-e}$). Claim~\ref{claimrelpsiil1i} also gives us
\begin{align*}
\alpha'_j &= \alpha'_{\omega_{t-1-e}}\\
&= \alpha'_{\omega_{\psi-\psi+t-1-e}}\\
&= \begin{cases}
\alpha'_{\omega_{\psi_{t-1-e-\psi}}} & \text{if } \psi \leq t-1-e\\
\alpha'_{\omega_{\psi_{t-1-e+t-\psi}}} & \text{if } \psi > t-1-e
\end{cases}\\
&= \begin{cases}
\alpha_{\omega_{t-1-(t-1-e-\psi)}} & \text{if } \psi \leq t-1-e\\
\alpha_{\omega_{t-1-(t-1-e+t-\psi)}} & \text{if } \psi > t-1-e
\end{cases}\\
&= \begin{cases}
\alpha_{\omega_{e+\psi}} & \text{if } \psi \leq t-1-e\\
\alpha_{\omega_{e-t+\psi}} & \text{if } \psi > t-1-e
\end{cases}\\
&= \alpha_{\omega_{\psi_e}}\\
&= \alpha_i.
\end{align*}
Note that $0 \leq \psi_e < t$ and hence $\psi_e = e + \psi$ if $\psi \leq t-1-e$ because $e + \psi \leq t-1$ and $\psi_e = e - t + \psi$ if $\psi > t-1-e$ because $e+\psi \geq t$.
\qed

By Claim~\ref{claimrelalphaialphaj} it also follows that the cycles of $\beta$ are disjoint: suppose $\beta$ contains $(\alpha_h,\alpha_i)$ and $(\alpha_j,\alpha_i)$ for some pairwise different $h,i,j \in [0,t-1]$. Then we have $\alpha_h = \alpha_i' = \alpha_j$ which yields $h = j$ contradicting $h \neq j$.

Now it suffices to show $l_\infty(\beta, \alpha^{t_1}) \leq 1$ and $l_\infty(\beta, \alpha^{t_2}) \leq 1$. By Claim~\ref{claimreldi1di2} we have for all $i \in [0,t-1]$ and all $k \in [1,2]$
\begin{displaymath}
\alpha_i^{\alpha^{t_k}} \in \{\alpha_j-1,\alpha_j,\alpha_j+1\}
\end{displaymath}
for some $j \in [0,t-1]$. Moreover we have $\alpha_i' = \alpha_j$. By Claim~\ref{claimrelalphaialphaj} we then have $\alpha_j' = \alpha_i$ and hence
\begin{displaymath}
\alpha_j^{\alpha^{t_k}} \in \{\alpha_i-1,\alpha_i,\alpha_i+1\}.
\end{displaymath}
Thus we either have $\beta_i = (\alpha_i, \alpha_j)$ and $\beta_j = \mathrm{id}$ or $\beta_i = \mathrm{id}$ and $\beta_j = (\alpha_i, \alpha_j)$ yielding a distance of at most $1$ if $i \neq j$. In the case
\begin{displaymath}
\alpha_i^{\alpha^{t_k}} \in \{\alpha_i-1,\alpha_i,\alpha_i+1\}
\end{displaymath}
we have that $\alpha_i$ is a fixed point in $\beta$ yielding a distance of at most $1$. Thus we finally obtain $l_\infty(\beta, \alpha^{t_1}) \leq 1$ and $l_\infty(\beta, \alpha^{t_2}) \leq 1$.
\end{proof}

\end{document}
